% Generated draft: requires mathematical and source review.
\documentclass[11pt]{article}
\usepackage[margin=25mm]{geometry}
\usepackage{fontspec}
\setmainfont{Latin Modern Roman}
\usepackage{amsmath,amssymb,mathtools,unicode-math}
\setmathfont{Latin Modern Math}
\usepackage{enumitem,needspace,xurl,hyperref}
\setlength{\parindent}{0pt}
\setlength{\parskip}{5pt}
\setlength{\emergencystretch}{3em}
\newcommand{\N}{\mathbb{N}}
\newcommand{\Z}{\mathbb{Z}}
\newcommand{\Q}{\mathbb{Q}}
\newcommand{\R}{\mathbb{R}}
\newcommand{\C}{\mathbb{C}}
\newcommand{\F}{\mathbb{F}}
\newcommand{\A}{\mathbb{A}}
\newcommand{\PP}{\mathbb{P}}
\newcommand{\E}{\mathbb{E}}
\newcommand{\eps}{\varepsilon}
\newcommand{\dd}{\,\mathrm{d}}
\begin{document}
\section*{UnsolvedMath ID: 30005563}
\textbf{Source problem code:} catalogue entry
\subsection{Definitions and mathematical statement}
The supplied record poses a collection of questions, rather than a single fully specified conjecture. Its hypothesis is that the vertex algebra \(V\) is not of CohFT-type, where CohFT means cohomological field theory. The excerpt does not specify the definition of CohFT-type, the admissible vertex algebras and module data, or the base parameter space. These hypotheses must be recovered before a definitive formulation is possible; being outside the stated CohFT-type class should not silently be replaced by another condition concerning cohomological field theories.

Let \(X\) denote a base parameter space intended by the source, and let \(\mathcal E\) be an associated sheaf of coinvariants. Whenever \(\mathcal E\) is locally free of finite rank \(r\), it is a vector bundle, with total Chern class
\[
c(\mathcal E)=1+c_1(\mathcal E)+\cdots+c_r(\mathcal E).
\]
Here the classes must be interpreted in a specified intersection or cohomology theory on \(X\). The source does not select that theory or its coefficients. If the chosen theory contains a specified tautological subring \(R^*(X)\), tautologicality of \(c_i(\mathcal E)\) means
\[
c_i(\mathcal E)\in R^i(X).
\]
Neither \(X\) nor this tautological subring is identified in the supplied excerpt. In a specified cycle theory, an effective class means a class represented by a cycle with nonnegative coefficients; the coefficient convention and any passage to numerical equivalence must also be fixed.

Subject to these unresolved choices, the source asks the following questions for coinvariants associated to \(V\) outside the CohFT-type setting.
\begin{enumerate}
\item Compute the Chern classes of the associated bundles. Determine when the higher Chern classes \(c_i(\mathcal E)\), for \(2\leq i\leq r\), are tautological.
\item Find sufficient conditions ensuring that the Chern classes \(c_i(\mathcal E)\), for \(1\leq i\leq r\), are effective. The original question requests sufficient conditions, without asserting that effectiveness always holds.
\item The excerpt attributes to the reference label \([\mathrm{CLTU23}]\) the existence of infinitely many extremal rays of an effective cone. For each ray \(\rho\) in that intended family, determine whether there is a bundle of coinvariants \(\mathcal E\) whose first Chern class spans \(\rho\). Explicitly, in the common real class space in which that cone is defined, the requested realization is
\[
\rho=\R_{\geq 0}\,c_1(\mathcal E),\qquad c_1(\mathcal E)\neq 0.
\]
A ray \(\rho\) of a cone \(C\) is extremal if \(u,v\in C\) and \(u+v\in\rho\) imply \(u,v\in\rho\). The space, cone, and family of rays attached to \([\mathrm{CLTU23}]\) are not identified in the supplied record. The distinction between a cone of effective classes and its closure is also unresolved.
\item Determine the degeneracy loci and their classes for the natural maps between sheaves of coinvariants referred to in the source. In the locally free case, for a specified map \(\phi:\mathcal E\to\mathcal F\) of bundles of ranks \(e\) and \(f\), and an integer \(0\leq k<\min(e,f)\), its rank degeneracy locus is
\[
D_k(\phi)=\{x\in X:\operatorname{rank}(\phi_x)\leq k\},
\]
with scheme structure given locally by the \((k+1)\)-minors. The relevant maps, rank thresholds, and meaning of their degeneracy classes remain unspecified. This bundle definition does not settle the source's broader question about maps of sheaves that may fail to be locally free.
\end{enumerate}
No local freeness, expected-codimension hypothesis, or extension of Chern classes to arbitrary coherent coinvariant sheaves is asserted by this formulation. Such hypotheses or constructions require separate specification.
\subsection{Short English statement}
For vertex algebras outside the source's CohFT-type setting, study the Chern classes of associated bundles of coinvariants, identify when higher classes are tautological, and give sufficient conditions for effectiveness. Ask whether the extremal effective rays attributed to \([\mathrm{CLTU23}]\) can be realized by first Chern classes, and study degeneracy loci and classes of natural maps between coinvariant sheaves. The base spaces, finiteness hypotheses, cycle theory, and cited family of rays are insufficiently specified for a definitive conjecture.
\subsection{Sources}
\begin{itemize}
\item Supplied local problem registry: permanent UnsolvedMath ID \(30005563\), title \emph{Chern and Degeneracy Classes of Coinvariant Bundles}, exact source problem code \texttt{catalogue entry}. The record lists proposed year 2023 and supplies no named proposer.
\item Primary report supplied by the registry: \emph{Recent Trends in Algebraic Geometry} (2023), Oberwolfach Reports, report 27/2023, pp. 1485--1542, DOI 10.4171/owr/2023/27. \url{https://doi.org/10.4171/owr/2023/27}. The registry describes the original excerpt as retained verbatim and labels its reconstructed statement \texttt{corrected\_verified}, with a supplied claim that it was checked against the official report. These are source-record qualifications, not independent verification performed for this draft.
\item Damiolini--Gibney--Krashen, \emph{Conformal blocks on smoothings via mode transition algebras}. Supplied literature link: \url{https://arxiv.org/abs/2307.03767}.
\item Damiolini--Gibney--Tarasca, \emph{Vertex algebras of CohFT-type}. Supplied literature link: \url{https://arxiv.org/abs/1910.01658}.
\item The original excerpt cites \([\mathrm{CLTU23}]\) for infinitely many extremal rays of an effective cone. No bibliographic identification or link for that reference is supplied, so its authors, target space, and precise conclusion remain unresolved.
\item Literature assessment dated exactly \texttt{2026-08-22}: the supplied curation labels the problem open and reports that primary work establishes coherence, factorization, vector-bundle cases, and tautological Chern classes under strong finiteness hypotheses, while broader questions for non-CohFT-type coinvariant sheaves and effective or extremal classes remain open. This is supplied dated curation derived from the OWR record, not an independently verified literature assessment. The record expressly requests independent verification of time-sensitive status and cited claims.
\end{itemize}
\end{document}
